\cgeentry{cge_mpsge}{Models in the form of MPSGE}

\cgesec{Description}

\cmd{cge_mpsge} (from version 0.44.0) writes a standalone model
(\manref{cge_model}) as MPSGE writes one \citep{rutherford1999}: sectors
whose activity levels make outputs from inputs, commodities whose prices
clear their markets, and consumers whose incomes come from endowments and
taxes. Each sector has a production block and each consumer a demand
block, whose records give quantities and prices at a reference point; the
functions are nested CES functions of the inputs and demands and CET
functions of the outputs, calibrated to those quantities and prices.
\cmd{cge_mpsge end} writes the model as user variables and equations
(\manref{cge_variable}, \manref{cge_equation}): the mixed complementarity
problem that MPSGE generates \citep{mathiesen1985, rutherford1995}, with zero
profit paired with each activity level, market clearing with each price and
income balance with each income. The model is then solved, reported and
exported like any standalone model.

\cgesec{Quick start}

A sector X making good PX from labour and capital, with a tax on labour,
and a consumer who owns the factors
\begin{stcode}
cge_model new
cge_parameter tx = 0.25
cge_mpsge sectors: X
cge_mpsge commodities: PX PL PK
cge_mpsge consumers: CONS
cge_mpsge prod:X s:0.5  o:PX q:100  i:PL q:40 p:1.25 a:CONS t:tx  i:PK q:50
cge_mpsge demand:CONS   d:PX q:100  e:PL q:40  e:PK q:50
cge_mpsge end, numeraire(PL)
cge_shock tx, value(0)
cge_solve, name(notax)
\end{stcode}
A model read from a GAMS program, between \cmd{$ontext} and \cmd{$offtext}
\begin{stcode}
cge_mpsge using "model.gms", numeraire(PF(lab))
\end{stcode}
List the definitions written; drop a model being read
\begin{stcode}
cge_mpsge, list generated
cge_mpsge, drop
\end{stcode}

\cgesec{Syntax}

A statement of the model, the model written, a model read from a file
\begin{stsyntax}
\cmdb{cge_mpsge} \argu{statement}\\
\cmdb{cge_mpsge end} [\texttt{,} \optarg{num}{eraire}{name[(labels)]} \optarg{st}{art}{levels}]\\
\cmdb{cge_mpsge using} \argu{filename} [\texttt{,} \optarg{num}{eraire}{name[(labels)]} \optarg{st}{art}{levels}]
\end{stsyntax}
List the models, or drop the model being read
\begin{stsyntax}
\cmdb{cge_mpsge} \texttt{,} \{\opt{}{list} [\opt{gen}{erated}]|\opt{}{drop}\}
\end{stsyntax}
A \argu{statement} is a section's header with what follows it, or the
fields of the block being read:
\begin{center}
\begin{tabular}{ll}
\hline
Statement & Content\\
\hline
\cmd{model:}\argu{name} & the model's name\\
\cmd{sectors:} \argu{names} & activity levels, as \cmd{X Y(s)}\\
\cmd{commodities:} \argu{names} & prices, as \cmd{PL P(g)}\\
\cmd{consumers:} \argu{names} & incomes, as \cmd{RA HH(h)}\\
\cmd{auxiliary:} \argu{names} & auxiliary variables, as \cmd{TAU U(r)}\\
\cmd{prod:}\argu{sector}[\cmd{$(}\argu{cond}\cmd{)}] [\argu{fields}] & a production block\\
\cmd{demand:}\argu{consumer}[\cmd{$(}\argu{cond}\cmd{)}] [\argu{fields}] & a demand block\\
\cmd{constraint:}\argu{aux}[\cmd{$(}\argu{cond}\cmd{)}] \argu{equation} & the constraint of an auxiliary variable\\
\cmd{report:} [\argu{reports}] & report variables\\
\argu{fields} & records and fields of the block being read\\
\hline
\end{tabular}
\end{center}
The fields are written \argu{key}\cmd{:}\argu{value}, as MPSGE writes them:
\begin{center}
\begin{tabular}{lp{0.62\textwidth}}
\hline
Field & Meaning\\
\hline
\cmd{s:}\argu{e} & elasticity of substitution at the top of the inputs or demands; default 0\\
\cmd{t:}\argu{e} & elasticity of transformation at the top of the outputs; default 0\\
\argu{nest}\cmd{:}\argu{e} & a nest within \cmd{s}; \argu{nest}\cmd{(}\argu{parent}\cmd{):}\argu{e} within another nest or \cmd{t}\\
\cmd{o:}, \cmd{i:} \argu{commodity} & an output or input of a production block\\
\cmd{d:}, \cmd{e:} \argu{commodity} & a final demand or endowment of a demand block\\
\cmd{q:}\argu{expression} & reference quantity of a record; default 1\\
\cmd{p:}\argu{expression} & reference price of a record; default 1\\
\cmd{a:}\argu{consumer} \cmd{t:}\argu{expression} & a tax at rate \cmd{t:}, paid to the agent \cmd{a:}\\
\cmd{n:}\argu{aux} [\cmd{m:}\argu{expression}] & an endogenous rate of the tax, \cmd{n:} times \cmd{m:} (default 1), added to \cmd{t:}\\
\cmd{r:}\argu{aux} & the multiplier of an endowment\\
\cmd{v:}\argu{name} & a report variable, with \cmd{o:} or \cmd{i:} and \cmd{prod:}, \cmd{d:} and \cmd{demand:}, or \cmd{w:}\argu{consumer}\\
\argu{nest}\cmd{:} & the nest of a record; default \cmd{s} for inputs and demands, \cmd{t} for outputs\\
\hline
\end{tabular}
\end{center}

\cgesec{Options}

\begin{optdesc}
\item[\optarg{num}{eraire}{name[(labels)]}] fixes the price of a commodity
or the income of a consumer at its value at the reference point, with
labels for an element of one with indexes, as \cmd{numeraire(PF(lab))}; its
equation is left out, as Walras's law makes it hold. The default is the
income of the consumer with the largest income at the reference point, as
in MPSGE.

\item[\optarg{st}{art}{levels}] (from version 0.45.0) gives the levels at
which sectors (default 1), commodities (default 1) and auxiliary variables
(default 0) start, as GAMS's \cmd{.l} before a solve, and at which the
incomes are evaluated at the reference point: \cmd{start(TAU:0.1 Z:0)}, or
\cmd{start(Y(agr):0)} for one element.

\item[\opt{}{list}] lists the model being read and the models written;
\opt{gen}{erated} lists the definitions written in place of the statements.
\opt{}{drop} drops the model being read. \manref{cge_equation}\cmd{, clear}
removes the models written with the other user variables and equations.
\end{optdesc}

\cgesec{Remarks and examples}

\cgesubsec{Statements}

Each \cmd{cge_mpsge} command holds a statement. A block may be written on
one line, with as many records as needed, or over several lines that each
continue the block before them. On a command line Stata reads
\cmd{$}\argu{name} as a global macro, so the sections are written without
MPSGE's dollar sign (\cmd{prod:} for \cmd{$prod:}) and conditions in
parentheses, \cmd{i:P(g)$(id0(g,s))}. \cmd{cge_mpsge using} reads a file as
MPSGE writes a model, dollar signs and conditions such as
\cmd{$id0(g,s)} included: from the line \cmd{$model:} to \cmd{$offtext},
or the whole file, so that it reads the model of a GAMS program between
\cmd{$ontext} and \cmd{$offtext}; lines that begin with an asterisk are
comments. The sets and parameters are those of the model in Stata, under
the same names (\manref{cge_set}, \manref{cge_parameter}); names are case
sensitive.

The fields are expressions in the parameters of the model. They are
evaluated at the parameters of each solve, as MPSGE generates the model
again before each solve, so that \manref{cge_shock} changes a tax, an
endowment or a technology. A block or record over indexes runs over their
sets: \cmd{prod:Y(s)} is a block for each element of \cmd{s}, and
\cmd{i:P(g) q:id0(g,s)} within it an input of each commodity \cmd{P(g)},
those with a reference quantity of zero left out, as in MPSGE. The
conditions, like those of \manref{cge_equation}, are evaluated at the base
unless \cmd{cge_options, conditions(solve)} evaluates them at each solve
(\manref{cge_options}).

\cgesubsec{Example 1: a tax on labour}

Two sectors make goods from labour and capital, with elasticities of 0.5
and 1.5; a third combines the goods into welfare, which the consumer buys
with the income of the factors and of a tax of 25 per cent on the labour of
X. The reference price of the taxed labour, 1.25, includes the tax, as MPSGE
takes the reference price of an input:
\stlog{mps_tax}
The data are a benchmark equilibrium, so that the equations hold at the
reference point. The definitions written are price indexes of the nests,
quantities of the records and the revenue of the tax, and three kinds of
equations:
\stlog{mps_list}
Without the tax, X expands and the price of its good falls; the
revenue, \cmd{X_tx2}, is 0:
\stlog{mps_solve}
The same model in GAMS/MPSGE gives the same equilibrium; the test suite
compares the two to $10^{-13}$.

\cgesubsec{Example 2: sectors over sets}

Three sectors use each other's goods as intermediate inputs, in fixed
proportions, and labour and capital in a nest whose elasticity varies by
sector; the government receives an output tax and a tax on labour in
manufacturing, and two households with elasticities of their own buy the
goods from the income of their endowments:
\stlog{mps_sets}
A tax of 50 per cent on labour in manufacturing shifts output toward the
other sectors and raises the government's income; \cmd{HH_u} is the
utility of each household relative to the reference point:
\stlog{mps_labtax}

\cgesubsec{What is written}

For sector $B$, consumer $H$, commodity $P$ and the $k$th record of a block:
\begin{center}
\begin{tabular}{lp{0.62\textwidth}}
\hline
Name & Content\\
\hline
$B$, $P$, $H$ & activity levels and prices, 1 at the reference point, and incomes at their reference values\\
$B$\cmd{_s}, $B$\cmd{_t}, $B$\cmd{_}\argu{nest} & price indexes of the nests, 1 at the reference point\\
$B$\cmd{_i}$k$, $B$\cmd{_o}$k$ & quantities of the inputs and outputs\\
$B$\cmd{_tx}$k$ & revenue of the tax on the record\\
$H$\cmd{_s}, $H$\cmd{_u}, $H$\cmd{_d}$k$ & price index and utility of the consumer, and its demands\\
$B$\cmd{_zp} & zero profit, paired with $B$\\
$P$\cmd{_mkt} & market clearing, supply less demand, paired with $P$\\
$H$\cmd{_inc} & income balance, paired with $H$\\
$A$\cmd{_con} & the constraint of auxiliary variable $A$, paired with it\\
\hline
\end{tabular}
\end{center}
The equations are box pairs with the lower bound 0 (\manref{cge_equation}):
a sector whose cost exceeds its revenue stays at 0, and a commodity in
excess supply has the price 0. A sector that does not pay at the reference
prices starts at 0 with \cmd{cge_bound} $B$\cmd{, level(0)}, as GAMS's
$B$\cmd{.l = 0} (\manref{cge_bound}). \cmd{cge_mpsge end} reports the
largest relative gap of the equations at the reference point, which is 0
for a benchmark equilibrium.

\cgesubsec{The numeraire}

Without \opt{num}{eraire()}, the income of the consumer with the largest
income at the reference point is fixed at that income. MPSGE fixes it at
its value at the start of each solve, so that prices and incomes may then
differ from those of GAMS by a common factor; quantities do not.

\cgesubsec{Example 3: auxiliary variables and constraints}

From version 0.45.0, an auxiliary variable is determined by its constraint,
\cmd{constraint:}\argu{aux} followed by an equation with \cmd{=g=},
\cmd{=e=} or \cmd{=l=}, which may go on over the statements that follow
it, to a semicolon or the next section; a consumer's name in a constraint
is its income. As in MPSGE the variable is at least 0 and starts at 0. It
may be the endogenous rate of a tax, \cmd{n:} (times \cmd{m:}, added to
\cmd{t:}), or the multiplier of an endowment, \cmd{r:}. Here the
government's tax on the output of X keeps its real revenue at
\cmd{grev}, and unemployment \cmd{U}, a negative endowment of labour
times \cmd{U}, keeps the real wage at least 1; \opt{st}{art()} starts the
tax at its rate at the reference point, so that the data are an
equilibrium there. Report variables give the labour of X and the welfare
index of the consumer:
\stlog{mps_aux}
With less capital the real wage would fall below its floor: a quarter of
the labour force is unemployed, and the tax rises to keep the
government's real revenue:
\stlog{mps_auxsolve}
The test suite compares this model with GAMS/MPSGE to $10^{-13}$.

\cgesubsec{Report variables}

\cmd{report:} defines variables from the model's: an output (\cmd{o:}) or
input (\cmd{i:}) of a production block, a demand (\cmd{d:}) of a demand
block, the sum over the block's records for the commodity, or the welfare
index of a consumer (\cmd{w:}), its utility $H$\cmd{_u}. A report variable
takes the indexes of its block and commodity that it names and sums over
the others: \cmd{v:LAB(s) i:PF(lab) prod:Y(s)} is the labour of each
sector, and \cmd{v:FAC(f) i:PF(f) prod:Y(man)} the factors of one.

\cgesubsec{What is not available}

Indexed nests stop \cmd{cge_mpsge} with an error.

\cgesec{Stored results}

\cmd{cge_mpsge end} and \cmd{cge_mpsge using} store the following in
\cmd{r()}:
\begin{storedres}
\resgroup{Scalars}
\res{r(N_sectors)} & activity levels (elements)\\
\res{r(N_commodities)} & prices (elements)\\
\res{r(N_consumers)} & incomes (elements)\\
\res{r(N_auxiliary)} & auxiliary variables (elements)\\
\res{r(N_reports)} & report variables\\
\res{r(N_blocks)} & blocks\\
\res{r(N_variables)} & user variables written\\
\res{r(N_equations)} & equations written, besides the definitions\\
\res{r(N_elements)} & equation elements written, with the definitions\\
\res{r(numeraire_value)} & value of the numeraire\\
\res{r(gap)} & largest relative gap of the equations at the reference point\\
\resgroup{Macros}
\res{r(name)} & the model's name (\cmd{mpsge} without \cmd{model:})\\
\res{r(numeraire)} & the numeraire\\
\end{storedres}
\cmd{cge_mpsge, list} stores \cmd{r(N_models)} and \cmd{r(models)}.

\cgesec{Methods and formulas}

In a nest with elasticity $e$, record $k$ with reference quantity $q_k$,
reference price $\bar p_k$, commodity price $P_k$ and tax rates summing to
$t_k$ has the value $v_k = \bar p_k q_k$ at reference prices and the
relative price $r_k = P_k (1 + t_k)/\bar p_k$ (an input or demand) or
$r_k = P_k (1 - t_k)/\bar p_k$ (an output). The price index of an input nest
is
\begin{equation*}
C = \Bigl[\frac{1}{V}\sum_k v_k\,r_k^{1-e}\Bigr]^{1/(1-e)}, \qquad V = \sum_k v_k,
\end{equation*}
over its records and inner nests, an inner nest with its own index in place
of $r_k$; the geometric mean at $e = 1$ and the weighted sum at $e = 0$. An
output nest has the exponent $1 + e$. A record's quantity per unit of
activity is $q_k (C/r_k)^{e}$, or $q_k (r_k/C)^{e}$ for an output, times the
factor of each nest above it. Zero profit of sector $B$ is
$V_s C_s - V_t R_t \ge 0 \perp B \ge 0$, with $C_s$ and $R_t$ the indexes of
the tops of its inputs and outputs. A consumer with income $M$ has utility
$U = M/(V_s C_s)$ and demands $q_k U$ times the factors of its nests. Market
clearing is supply, outputs and endowments, less demand, inputs and final
demands, $\ge 0 \perp P \ge 0$, and income balance $M = \sum P e + \sum
t P x \perp M \ge 0$, the value of the endowments and the revenue of the
taxes the consumer receives \citep{mathiesen1985, rutherford1999}. A tax
rate is $t + m n$, with $n$ an auxiliary variable; an endowment $q r$. The
constraint $f \ge g$ (or $f = g$) of an auxiliary variable $A$ is
$f - g \ge 0 \perp A \ge 0$, and $f \le g$ is $g - f \ge 0 \perp A \ge 0$. An
elasticity given by an expression is written with both forms, under the
conditions \cmd{$(}$e \ne 1$\cmd{)} and \cmd{$(}$e = 1$\cmd{)}.

\cgesec{References}

\citet{mathiesen1985}; \citet{rutherford1995}; \citet{rutherford1999}.

\cgesec{Also see}

\begin{alsosee}
\seeentry{cge_model}{A standalone model, written whole in the user language}
\seeentry{cge_variable}{Add a user variable}
\seeentry{cge_equation}{Add a user equation}
\seeentry{cge_nest}{Declare a nested CES or CET block}
\seeentry{cge_solve}{Solve the model}
\end{alsosee}
